Software systems in safety-critical domains such as medical imaging and
industrial automation evolve rapidly, and their runtime behaviour quickly
outgrows manual documentation.
\sysml Interaction Models capture scenario-specific component interactions,
but existing reconstruction techniques are each incomplete: static analysis
is scenario-agnostic, while execution traces are inherently partial.

This thesis presents a hybrid pipeline that synthesises scenario-specific
Interaction Models for event-driven \Cpp\ systems by combining a static call
graph, stored in a \neo property graph, with timestamped runtime traces
collected by a separate pipeline. The two evidence sources are kept in
separate artefacts and combined as a set union of their interaction edges;
scenario-scoped Cypher queries then isolate the interactions relevant to a
given scenario.
Implemented as eleven composable, natural-language-invocable agent skills,
the pipeline is evaluated on \emph{CPSCore}, an open-source event-driven
framework, across four scenarios (three happy-path, one constructed)
reconstructed under five conditions ranging from an \ac{LLM}-only baseline
to the full pipeline.

Three hypotheses are confirmed.
\textbf{H1:} The full pipeline achieves the highest reconstruction accuracy
of any condition tested (macro-F1~$=$~0.964, with perfect recall on every
scenario). Combining the two sources is strictly additive where static and
dynamic analysis each miss a different interaction and redundant where the
static call graph is already complete.
\textbf{H2:} Code-graph-guided instrumentation eliminates all measured noise
(50\% fewer nodes, 70\% fewer edges) without losing scenario coverage. It is also
confirmed by a blinded six-rater expert panel rating guided output higher in 47 of 54
comparisons.
\textbf{H3:} The reconstructed diagram alone achieves perfect component
recall, interaction recall and completeness on every scenario, either
matching or exceeding all other evidence combinations.
Together, the three findings converge on one pattern: additional evidence
improves reconstruction or comprehension specifically when the primary
evidence source is missing something.
As an exploratory complement, a practitioner session on a second, larger
industrial codebase found the underlying static-extraction and query
approach transferred well. Structural and dependency diagrams were rated
consistently useful and sequence-diagram message labelling was identified
as the clearest area for improvement.

The main contributions are as follows.
(i)~A static property graph plus a separate runtime-trace pipeline whose
interaction edge sets are compared while remembering which evidence each
interaction came from.
(ii)~A scenario-scoping method, using plain scenario-scoped Cypher queries,
that trims the call graph down to just the interactions relevant to one
scenario.
(iii)~A way to recover publish-subscribe wiring that ordinary call-graph
analysis cannot see.
(iv)~A small, reversible tool (\texttt{csp\_matcher}) that inserts and
removes trace probes automatically.
(v)~A set of eleven natural-language agent skills that run the whole
pipeline, from source code to \sysml diagram. Its designed to be codebase agnostic.
(vi)~A clear, evidence-based account of when combining static and dynamic
analysis actually helps and when it is redundant.
